% Options for packages loaded elsewhere
\PassOptionsToPackage{unicode}{hyperref}
\PassOptionsToPackage{hyphens}{url}
\documentclass[
]{article}
\usepackage{xcolor}
\usepackage[margin=1in]{geometry}
\usepackage{amsmath,amssymb}
\setcounter{secnumdepth}{-\maxdimen} % remove section numbering
\usepackage{iftex}
\ifPDFTeX
  \usepackage[T1]{fontenc}
  \usepackage[utf8]{inputenc}
  \usepackage{textcomp} % provide euro and other symbols
\else % if luatex or xetex
  \usepackage{unicode-math} % this also loads fontspec
  \defaultfontfeatures{Scale=MatchLowercase}
  \defaultfontfeatures[\rmfamily]{Ligatures=TeX,Scale=1}
\fi
\usepackage{lmodern}
\ifPDFTeX\else
  % xetex/luatex font selection
\fi
% Use upquote if available, for straight quotes in verbatim environments
\IfFileExists{upquote.sty}{\usepackage{upquote}}{}
\IfFileExists{microtype.sty}{% use microtype if available
  \usepackage[]{microtype}
  \UseMicrotypeSet[protrusion]{basicmath} % disable protrusion for tt fonts
}{}
\makeatletter
\@ifundefined{KOMAClassName}{% if non-KOMA class
  \IfFileExists{parskip.sty}{%
    \usepackage{parskip}
  }{% else
    \setlength{\parindent}{0pt}
    \setlength{\parskip}{6pt plus 2pt minus 1pt}}
}{% if KOMA class
  \KOMAoptions{parskip=half}}
\makeatother
\usepackage{color}
\usepackage{fancyvrb}
\newcommand{\VerbBar}{|}
\newcommand{\VERB}{\Verb[commandchars=\\\{\}]}
\DefineVerbatimEnvironment{Highlighting}{Verbatim}{commandchars=\\\{\}}
% Add ',fontsize=\small' for more characters per line
\usepackage{framed}
\definecolor{shadecolor}{RGB}{248,248,248}
\newenvironment{Shaded}{\begin{snugshade}}{\end{snugshade}}
\newcommand{\AlertTok}[1]{\textcolor[rgb]{0.94,0.16,0.16}{#1}}
\newcommand{\AnnotationTok}[1]{\textcolor[rgb]{0.56,0.35,0.01}{\textbf{\textit{#1}}}}
\newcommand{\AttributeTok}[1]{\textcolor[rgb]{0.13,0.29,0.53}{#1}}
\newcommand{\BaseNTok}[1]{\textcolor[rgb]{0.00,0.00,0.81}{#1}}
\newcommand{\BuiltInTok}[1]{#1}
\newcommand{\CharTok}[1]{\textcolor[rgb]{0.31,0.60,0.02}{#1}}
\newcommand{\CommentTok}[1]{\textcolor[rgb]{0.56,0.35,0.01}{\textit{#1}}}
\newcommand{\CommentVarTok}[1]{\textcolor[rgb]{0.56,0.35,0.01}{\textbf{\textit{#1}}}}
\newcommand{\ConstantTok}[1]{\textcolor[rgb]{0.56,0.35,0.01}{#1}}
\newcommand{\ControlFlowTok}[1]{\textcolor[rgb]{0.13,0.29,0.53}{\textbf{#1}}}
\newcommand{\DataTypeTok}[1]{\textcolor[rgb]{0.13,0.29,0.53}{#1}}
\newcommand{\DecValTok}[1]{\textcolor[rgb]{0.00,0.00,0.81}{#1}}
\newcommand{\DocumentationTok}[1]{\textcolor[rgb]{0.56,0.35,0.01}{\textbf{\textit{#1}}}}
\newcommand{\ErrorTok}[1]{\textcolor[rgb]{0.64,0.00,0.00}{\textbf{#1}}}
\newcommand{\ExtensionTok}[1]{#1}
\newcommand{\FloatTok}[1]{\textcolor[rgb]{0.00,0.00,0.81}{#1}}
\newcommand{\FunctionTok}[1]{\textcolor[rgb]{0.13,0.29,0.53}{\textbf{#1}}}
\newcommand{\ImportTok}[1]{#1}
\newcommand{\InformationTok}[1]{\textcolor[rgb]{0.56,0.35,0.01}{\textbf{\textit{#1}}}}
\newcommand{\KeywordTok}[1]{\textcolor[rgb]{0.13,0.29,0.53}{\textbf{#1}}}
\newcommand{\NormalTok}[1]{#1}
\newcommand{\OperatorTok}[1]{\textcolor[rgb]{0.81,0.36,0.00}{\textbf{#1}}}
\newcommand{\OtherTok}[1]{\textcolor[rgb]{0.56,0.35,0.01}{#1}}
\newcommand{\PreprocessorTok}[1]{\textcolor[rgb]{0.56,0.35,0.01}{\textit{#1}}}
\newcommand{\RegionMarkerTok}[1]{#1}
\newcommand{\SpecialCharTok}[1]{\textcolor[rgb]{0.81,0.36,0.00}{\textbf{#1}}}
\newcommand{\SpecialStringTok}[1]{\textcolor[rgb]{0.31,0.60,0.02}{#1}}
\newcommand{\StringTok}[1]{\textcolor[rgb]{0.31,0.60,0.02}{#1}}
\newcommand{\VariableTok}[1]{\textcolor[rgb]{0.00,0.00,0.00}{#1}}
\newcommand{\VerbatimStringTok}[1]{\textcolor[rgb]{0.31,0.60,0.02}{#1}}
\newcommand{\WarningTok}[1]{\textcolor[rgb]{0.56,0.35,0.01}{\textbf{\textit{#1}}}}
\usepackage{graphicx}
\makeatletter
\newsavebox\pandoc@box
\newcommand*\pandocbounded[1]{% scales image to fit in text height/width
  \sbox\pandoc@box{#1}%
  \Gscale@div\@tempa{\textheight}{\dimexpr\ht\pandoc@box+\dp\pandoc@box\relax}%
  \Gscale@div\@tempb{\linewidth}{\wd\pandoc@box}%
  \ifdim\@tempb\p@<\@tempa\p@\let\@tempa\@tempb\fi% select the smaller of both
  \ifdim\@tempa\p@<\p@\scalebox{\@tempa}{\usebox\pandoc@box}%
  \else\usebox{\pandoc@box}%
  \fi%
}
% Set default figure placement to htbp
\def\fps@figure{htbp}
\makeatother
\setlength{\emergencystretch}{3em} % prevent overfull lines
\providecommand{\tightlist}{%
  \setlength{\itemsep}{0pt}\setlength{\parskip}{0pt}}
\usepackage{bookmark}
\IfFileExists{xurl.sty}{\usepackage{xurl}}{} % add URL line breaks if available
\urlstyle{same}
\hypersetup{
  pdftitle={analisis\_IPC},
  pdfauthor={Gloria Disney Restrepo Ruiz},
  hidelinks,
  pdfcreator={LaTeX via pandoc}}

\title{analisis\_IPC}
\author{Gloria Disney Restrepo Ruiz}
\date{2026-08-29}

\begin{document}
\maketitle

\begin{Shaded}
\begin{Highlighting}[]
\FunctionTok{getwd}\NormalTok{()}
\end{Highlighting}
\end{Shaded}

\begin{verbatim}
## [1] "C:/Users/disne/OneDrive/Desktop/ipc"
\end{verbatim}

\begin{Shaded}
\begin{Highlighting}[]
\NormalTok{ruta\_crudos }\OtherTok{\textless{}{-}} \FunctionTok{here}\NormalTok{(}\StringTok{"datos"}\NormalTok{, }\StringTok{"crudos"}\NormalTok{)}

\NormalTok{archivos }\OtherTok{\textless{}{-}} \FunctionTok{list.files}\NormalTok{(ruta\_crudos, }\AttributeTok{pattern =} \StringTok{"\^{}anex{-}IPC{-}.*2026}\SpecialCharTok{\textbackslash{}\textbackslash{}}\StringTok{.xlsx$"}\NormalTok{, }\AttributeTok{full.names =} \ConstantTok{TRUE}\NormalTok{)}
\NormalTok{archivos   }\CommentTok{\# si esto sale character(0), revisa que los 7 .xlsx estén en datos/crudos/}
\end{Highlighting}
\end{Shaded}

\begin{verbatim}
## [1] "C:/Users/disne/OneDrive/Desktop/ipc/datos/crudos/anex-IPC-abr2026.xlsx"
## [2] "C:/Users/disne/OneDrive/Desktop/ipc/datos/crudos/anex-IPC-ene2026.xlsx"
## [3] "C:/Users/disne/OneDrive/Desktop/ipc/datos/crudos/anex-IPC-feb2026.xlsx"
## [4] "C:/Users/disne/OneDrive/Desktop/ipc/datos/crudos/anex-IPC-jul2026.xlsx"
## [5] "C:/Users/disne/OneDrive/Desktop/ipc/datos/crudos/anex-IPC-jun2026.xlsx"
## [6] "C:/Users/disne/OneDrive/Desktop/ipc/datos/crudos/anex-IPC-mar2026.xlsx"
## [7] "C:/Users/disne/OneDrive/Desktop/ipc/datos/crudos/anex-IPC-may2026.xlsx"
\end{verbatim}

\begin{Shaded}
\begin{Highlighting}[]
\CommentTok{\# Tabla de referencia para traducir la abreviatura del nombre de archivo a mes}
\NormalTok{meses\_lookup }\OtherTok{\textless{}{-}} \FunctionTok{tibble}\NormalTok{(}
  \AttributeTok{abrev      =} \FunctionTok{c}\NormalTok{(}\StringTok{"ene"}\NormalTok{,}\StringTok{"feb"}\NormalTok{,}\StringTok{"mar"}\NormalTok{,}\StringTok{"abr"}\NormalTok{,}\StringTok{"may"}\NormalTok{,}\StringTok{"jun"}\NormalTok{,}\StringTok{"jul"}\NormalTok{),}
  \AttributeTok{mes\_num    =} \DecValTok{1}\SpecialCharTok{:}\DecValTok{7}\NormalTok{,}
  \AttributeTok{mes\_nombre =} \FunctionTok{c}\NormalTok{(}\StringTok{"Enero"}\NormalTok{,}\StringTok{"Febrero"}\NormalTok{,}\StringTok{"Marzo"}\NormalTok{,}\StringTok{"Abril"}\NormalTok{,}\StringTok{"Mayo"}\NormalTok{,}\StringTok{"Junio"}\NormalTok{,}\StringTok{"Julio"}\NormalTok{)}
\NormalTok{)}

\NormalTok{extraer\_abrev }\OtherTok{\textless{}{-}} \ControlFlowTok{function}\NormalTok{(nombre\_archivo) \{}
  \FunctionTok{str\_extract}\NormalTok{(}\FunctionTok{basename}\NormalTok{(nombre\_archivo), }\StringTok{"(?\textless{}=anex{-}IPC{-})[a{-}z]\{3\}(?=2026)"}\NormalTok{)}
\NormalTok{\}}
\end{Highlighting}
\end{Shaded}

\begin{Shaded}
\begin{Highlighting}[]
\NormalTok{leer\_total\_mes }\OtherTok{\textless{}{-}} \ControlFlowTok{function}\NormalTok{(archivo) \{}
\NormalTok{  abrev }\OtherTok{\textless{}{-}} \FunctionTok{extraer\_abrev}\NormalTok{(archivo)}
\NormalTok{  datos }\OtherTok{\textless{}{-}} \FunctionTok{read\_excel}\NormalTok{(}
\NormalTok{    archivo, }\AttributeTok{sheet =} \StringTok{"1"}\NormalTok{, }\AttributeTok{skip =} \DecValTok{9}\NormalTok{, }\AttributeTok{n\_max =} \DecValTok{10}\NormalTok{,}
    \AttributeTok{col\_names =} \FunctionTok{c}\NormalTok{(}\StringTok{"anio"}\NormalTok{, }\StringTok{"var\_mensual"}\NormalTok{, }\StringTok{"var\_anio\_corrido"}\NormalTok{, }\StringTok{"var\_anual"}\NormalTok{)}
\NormalTok{  )}
\NormalTok{  datos }\SpecialCharTok{\%\textgreater{}\%}
    \FunctionTok{filter}\NormalTok{(anio }\SpecialCharTok{==} \DecValTok{2026}\NormalTok{) }\SpecialCharTok{\%\textgreater{}\%}
    \FunctionTok{mutate}\NormalTok{(}\AttributeTok{mes\_abrev =}\NormalTok{ abrev)}
\NormalTok{\}}

\NormalTok{serie\_total }\OtherTok{\textless{}{-}} \FunctionTok{map\_dfr}\NormalTok{(archivos, leer\_total\_mes) }\SpecialCharTok{\%\textgreater{}\%}
  \FunctionTok{left\_join}\NormalTok{(meses\_lookup, }\AttributeTok{by =} \FunctionTok{c}\NormalTok{(}\StringTok{"mes\_abrev"} \OtherTok{=} \StringTok{"abrev"}\NormalTok{)) }\SpecialCharTok{\%\textgreater{}\%}
  \FunctionTok{mutate}\NormalTok{(}\AttributeTok{fecha =} \FunctionTok{as.Date}\NormalTok{(}\FunctionTok{paste}\NormalTok{(}\StringTok{"2026"}\NormalTok{, mes\_num, }\StringTok{"01"}\NormalTok{, }\AttributeTok{sep =} \StringTok{"{-}"}\NormalTok{))) }\SpecialCharTok{\%\textgreater{}\%}
  \FunctionTok{arrange}\NormalTok{(fecha) }\SpecialCharTok{\%\textgreater{}\%}
  \FunctionTok{select}\NormalTok{(fecha, }\AttributeTok{mes =}\NormalTok{ mes\_nombre, var\_mensual, var\_anio\_corrido, var\_anual)}

\NormalTok{serie\_total}
\end{Highlighting}
\end{Shaded}

\begin{verbatim}
## # A tibble: 7 x 5
##   fecha      mes     var_mensual var_anio_corrido var_anual
##   <date>     <chr>         <dbl>            <dbl>     <dbl>
## 1 2026-01-01 Enero          1.18             1.18      5.35
## 2 2026-02-01 Febrero        1.08             2.27      5.29
## 3 2026-03-01 Marzo          0.78             3.07      5.56
## 4 2026-04-01 Abril          0.78             3.87      5.68
## 5 2026-05-01 Mayo           0.47             4.36      5.84
## 6 2026-06-01 Junio          0.39             4.77      6.14
## 7 2026-07-01 Julio          0.17             4.94      6.03
\end{verbatim}

\begin{Shaded}
\begin{Highlighting}[]
\NormalTok{leer\_divisiones\_mes }\OtherTok{\textless{}{-}} \ControlFlowTok{function}\NormalTok{(archivo) \{}
\NormalTok{  abrev }\OtherTok{\textless{}{-}} \FunctionTok{extraer\_abrev}\NormalTok{(archivo)}
  \CommentTok{\# La hoja "2" trae 11 columnas de datos (Variación, Contribución y}
  \CommentTok{\# Participación), pero solo necesitamos las 5 primeras (A:E). Por eso usamos}
  \CommentTok{\# "range" en vez de "skip", para quedarnos únicamente con esas columnas.}
\NormalTok{  datos }\OtherTok{\textless{}{-}} \FunctionTok{read\_excel}\NormalTok{(}
\NormalTok{    archivo, }\AttributeTok{sheet =} \StringTok{"2"}\NormalTok{,}
    \AttributeTok{range =}\NormalTok{ readxl}\SpecialCharTok{::}\FunctionTok{cell\_limits}\NormalTok{(}\FunctionTok{c}\NormalTok{(}\DecValTok{10}\NormalTok{, }\DecValTok{1}\NormalTok{), }\FunctionTok{c}\NormalTok{(}\DecValTok{21}\NormalTok{, }\DecValTok{5}\NormalTok{)),}
    \AttributeTok{col\_names =} \FunctionTok{c}\NormalTok{(}\StringTok{"division"}\NormalTok{, }\StringTok{"ponderacion"}\NormalTok{, }\StringTok{"var\_mensual"}\NormalTok{, }\StringTok{"var\_anio\_corrido"}\NormalTok{, }\StringTok{"var\_anual"}\NormalTok{)}
\NormalTok{  )}
\NormalTok{  datos }\SpecialCharTok{\%\textgreater{}\%} \FunctionTok{mutate}\NormalTok{(}\AttributeTok{mes\_abrev =}\NormalTok{ abrev)}
\NormalTok{\}}
\end{Highlighting}
\end{Shaded}

\begin{Shaded}
\begin{Highlighting}[]
\CommentTok{\# Este chunk faltaba: aquí sí se USA la función de arriba para construir}
\CommentTok{\# el objeto "divisiones" (antes el archivo pasaba directo a usarlo sin crearlo).}
\NormalTok{divisiones }\OtherTok{\textless{}{-}} \FunctionTok{map\_dfr}\NormalTok{(archivos, leer\_divisiones\_mes) }\SpecialCharTok{\%\textgreater{}\%}
  \FunctionTok{left\_join}\NormalTok{(meses\_lookup, }\AttributeTok{by =} \FunctionTok{c}\NormalTok{(}\StringTok{"mes\_abrev"} \OtherTok{=} \StringTok{"abrev"}\NormalTok{)) }\SpecialCharTok{\%\textgreater{}\%}
  \FunctionTok{mutate}\NormalTok{(}
    \AttributeTok{division =} \FunctionTok{str\_squish}\NormalTok{(division),}
    \AttributeTok{fecha =} \FunctionTok{as.Date}\NormalTok{(}\FunctionTok{paste}\NormalTok{(}\StringTok{"2026"}\NormalTok{, mes\_num, }\StringTok{"01"}\NormalTok{, }\AttributeTok{sep =} \StringTok{"{-}"}\NormalTok{))}
\NormalTok{  ) }\SpecialCharTok{\%\textgreater{}\%}
  \FunctionTok{arrange}\NormalTok{(fecha, division) }\SpecialCharTok{\%\textgreater{}\%}
  \FunctionTok{select}\NormalTok{(fecha, }\AttributeTok{mes =}\NormalTok{ mes\_nombre, division, ponderacion, var\_mensual, var\_anio\_corrido, var\_anual)}

\NormalTok{divisiones}
\end{Highlighting}
\end{Shaded}

\begin{verbatim}
## # A tibble: 84 x 7
##    fecha      mes   division  ponderacion var_mensual var_anio_corrido var_anual
##    <date>     <chr> <chr>           <dbl>       <dbl>            <dbl>     <dbl>
##  1 2026-01-01 Enero Alimento~       15.0         1.66             1.66      5.11
##  2 2026-01-01 Enero Alojamie~       33.1         0.23             0.23      4.59
##  3 2026-01-01 Enero Bebidas ~        1.7         1.79             1.79      7.58
##  4 2026-01-01 Enero Bienes Y~        5.36        0.79             0.79      5.23
##  5 2026-01-01 Enero Educación        4.41        0                0         7.36
##  6 2026-01-01 Enero Informac~        4.33        0.06             0.06      1.48
##  7 2026-01-01 Enero Muebles,~        4.19        1.5              1.5       4.12
##  8 2026-01-01 Enero Prendas ~        3.98        0.18             0.18      2.45
##  9 2026-01-01 Enero Recreaci~        3.79        0.5              0.5       2.54
## 10 2026-01-01 Enero Restaura~        9.43        2.94             2.94      9.01
## # i 74 more rows
\end{verbatim}

\begin{Shaded}
\begin{Highlighting}[]
\CommentTok{\# Verificación: 7 meses en la serie total, 12 divisiones x 7 meses en el detalle}
\FunctionTok{nrow}\NormalTok{(serie\_total)   }\CommentTok{\# esperado: 7}
\end{Highlighting}
\end{Shaded}

\begin{verbatim}
## [1] 7
\end{verbatim}

\begin{Shaded}
\begin{Highlighting}[]
\FunctionTok{nrow}\NormalTok{(divisiones)    }\CommentTok{\# esperado: 84}
\end{Highlighting}
\end{Shaded}

\begin{verbatim}
## [1] 84
\end{verbatim}

\begin{Shaded}
\begin{Highlighting}[]
\CommentTok{\# Verificación de valores faltantes}
\FunctionTok{colSums}\NormalTok{(}\FunctionTok{is.na}\NormalTok{(serie\_total))}
\end{Highlighting}
\end{Shaded}

\begin{verbatim}
##            fecha              mes      var_mensual var_anio_corrido 
##                0                0                0                0 
##        var_anual 
##                0
\end{verbatim}

\begin{Shaded}
\begin{Highlighting}[]
\FunctionTok{colSums}\NormalTok{(}\FunctionTok{is.na}\NormalTok{(divisiones))}
\end{Highlighting}
\end{Shaded}

\begin{verbatim}
##            fecha              mes         division      ponderacion 
##                0                0                0                0 
##      var_mensual var_anio_corrido        var_anual 
##                0                0                0
\end{verbatim}

\begin{Shaded}
\begin{Highlighting}[]
\CommentTok{\# Guardar datos procesados}
\FunctionTok{dir.create}\NormalTok{(}\FunctionTok{here}\NormalTok{(}\StringTok{"datos"}\NormalTok{, }\StringTok{"procesados"}\NormalTok{), }\AttributeTok{showWarnings =} \ConstantTok{FALSE}\NormalTok{)}
\FunctionTok{write.csv}\NormalTok{(serie\_total, }\FunctionTok{here}\NormalTok{(}\StringTok{"datos"}\NormalTok{, }\StringTok{"procesados"}\NormalTok{, }\StringTok{"serie\_total\_ipc.csv"}\NormalTok{), }\AttributeTok{row.names =} \ConstantTok{FALSE}\NormalTok{)}
\FunctionTok{write.csv}\NormalTok{(divisiones, }\FunctionTok{here}\NormalTok{(}\StringTok{"datos"}\NormalTok{, }\StringTok{"procesados"}\NormalTok{, }\StringTok{"divisiones\_ipc.csv"}\NormalTok{), }\AttributeTok{row.names =} \ConstantTok{FALSE}\NormalTok{)}
\end{Highlighting}
\end{Shaded}

\begin{Shaded}
\begin{Highlighting}[]
\FunctionTok{summary}\NormalTok{(serie\_total}\SpecialCharTok{$}\NormalTok{var\_mensual)}
\end{Highlighting}
\end{Shaded}

\begin{verbatim}
##    Min. 1st Qu.  Median    Mean 3rd Qu.    Max. 
##  0.1700  0.4300  0.7800  0.6929  0.9300  1.1800
\end{verbatim}

\begin{Shaded}
\begin{Highlighting}[]
\FunctionTok{summary}\NormalTok{(serie\_total}\SpecialCharTok{$}\NormalTok{var\_anual)}
\end{Highlighting}
\end{Shaded}

\begin{verbatim}
##    Min. 1st Qu.  Median    Mean 3rd Qu.    Max. 
##   5.290   5.455   5.680   5.699   5.935   6.140
\end{verbatim}

\begin{Shaded}
\begin{Highlighting}[]
\CommentTok{\# Gráfico 1: evolución de la variación mensual y anual del IPC (ene{-}jul 2026)}
\NormalTok{serie\_larga }\OtherTok{\textless{}{-}}\NormalTok{ serie\_total }\SpecialCharTok{\%\textgreater{}\%}
  \FunctionTok{select}\NormalTok{(fecha, mes, var\_mensual, var\_anual) }\SpecialCharTok{\%\textgreater{}\%}
  \FunctionTok{pivot\_longer}\NormalTok{(}\AttributeTok{cols =} \FunctionTok{c}\NormalTok{(var\_mensual, var\_anual), }\AttributeTok{names\_to =} \StringTok{"tipo"}\NormalTok{, }\AttributeTok{values\_to =} \StringTok{"valor"}\NormalTok{) }\SpecialCharTok{\%\textgreater{}\%}
  \FunctionTok{mutate}\NormalTok{(}\AttributeTok{tipo =} \FunctionTok{recode}\NormalTok{(tipo,}
    \AttributeTok{var\_mensual =} \StringTok{"Variación mensual"}\NormalTok{,}
    \AttributeTok{var\_anual   =} \StringTok{"Variación anual"}
\NormalTok{  ))}

\NormalTok{grafico\_tendencia }\OtherTok{\textless{}{-}} \FunctionTok{ggplot}\NormalTok{(serie\_larga, }\FunctionTok{aes}\NormalTok{(}\AttributeTok{x =}\NormalTok{ fecha, }\AttributeTok{y =}\NormalTok{ valor, }\AttributeTok{color =}\NormalTok{ tipo)) }\SpecialCharTok{+}
  \FunctionTok{geom\_line}\NormalTok{(}\AttributeTok{linewidth =} \FloatTok{1.1}\NormalTok{) }\SpecialCharTok{+}
  \FunctionTok{geom\_point}\NormalTok{(}\AttributeTok{size =} \DecValTok{2}\NormalTok{) }\SpecialCharTok{+}
  \FunctionTok{scale\_x\_date}\NormalTok{(}\AttributeTok{date\_labels =} \StringTok{"\%b"}\NormalTok{, }\AttributeTok{date\_breaks =} \StringTok{"1 month"}\NormalTok{) }\SpecialCharTok{+}
  \FunctionTok{labs}\NormalTok{(}
    \AttributeTok{title =} \StringTok{"Variación del IPC en Colombia, enero{-}julio 2026"}\NormalTok{,}
    \AttributeTok{subtitle =} \StringTok{"Variación mensual vs. variación anual (interanual)"}\NormalTok{,}
    \AttributeTok{x =} \ConstantTok{NULL}\NormalTok{, }\AttributeTok{y =} \StringTok{"Variación (\%)"}\NormalTok{, }\AttributeTok{color =} \ConstantTok{NULL}\NormalTok{,}
    \AttributeTok{caption =} \StringTok{"Fuente: DANE, boletines técnicos mensuales del IPC (anexos)"}
\NormalTok{  ) }\SpecialCharTok{+}
  \FunctionTok{theme\_minimal}\NormalTok{(}\AttributeTok{base\_size =} \DecValTok{13}\NormalTok{) }\SpecialCharTok{+}
  \FunctionTok{theme}\NormalTok{(}\AttributeTok{legend.position =} \StringTok{"bottom"}\NormalTok{)}

\NormalTok{grafico\_tendencia}
\end{Highlighting}
\end{Shaded}

\pandocbounded{\includegraphics[keepaspectratio]{analisis_ipc_files/analisis_ipc_files/figure-latex/unnamed-chunk-8-1.pdf}}

\begin{Shaded}
\begin{Highlighting}[]
\FunctionTok{dir.create}\NormalTok{(}\FunctionTok{here}\NormalTok{(}\StringTok{"salidas"}\NormalTok{), }\AttributeTok{showWarnings =} \ConstantTok{FALSE}\NormalTok{)}
\FunctionTok{ggsave}\NormalTok{(}\FunctionTok{here}\NormalTok{(}\StringTok{"salidas"}\NormalTok{, }\StringTok{"01\_tendencia\_ipc.png"}\NormalTok{), grafico\_tendencia, }\AttributeTok{width =} \DecValTok{8}\NormalTok{, }\AttributeTok{height =} \DecValTok{5}\NormalTok{, }\AttributeTok{dpi =} \DecValTok{150}\NormalTok{)}
\end{Highlighting}
\end{Shaded}

\begin{Shaded}
\begin{Highlighting}[]
\CommentTok{\# Gráfico 2: variación mensual por división de gasto, julio 2026 (mes más reciente)}
\NormalTok{division\_julio }\OtherTok{\textless{}{-}}\NormalTok{ divisiones }\SpecialCharTok{\%\textgreater{}\%} \FunctionTok{filter}\NormalTok{(mes }\SpecialCharTok{==} \StringTok{"Julio"}\NormalTok{)}

\NormalTok{grafico\_divisiones }\OtherTok{\textless{}{-}} \FunctionTok{ggplot}\NormalTok{(division\_julio, }\FunctionTok{aes}\NormalTok{(}\AttributeTok{x =} \FunctionTok{reorder}\NormalTok{(division, var\_mensual), }\AttributeTok{y =}\NormalTok{ var\_mensual)) }\SpecialCharTok{+}
  \FunctionTok{geom\_col}\NormalTok{(}\AttributeTok{fill =} \StringTok{"\#2E86AB"}\NormalTok{) }\SpecialCharTok{+}
  \FunctionTok{coord\_flip}\NormalTok{() }\SpecialCharTok{+}
  \FunctionTok{labs}\NormalTok{(}
    \AttributeTok{title =} \StringTok{"Variación mensual del IPC por división de gasto"}\NormalTok{,}
    \AttributeTok{subtitle =} \StringTok{"Julio de 2026"}\NormalTok{,}
    \AttributeTok{x =} \ConstantTok{NULL}\NormalTok{, }\AttributeTok{y =} \StringTok{"Variación mensual (\%)"}\NormalTok{,}
    \AttributeTok{caption =} \StringTok{"Fuente: DANE, boletín técnico IPC julio 2026 (anexo)"}
\NormalTok{  ) }\SpecialCharTok{+}
  \FunctionTok{theme\_minimal}\NormalTok{(}\AttributeTok{base\_size =} \DecValTok{12}\NormalTok{)}

\NormalTok{grafico\_divisiones}
\end{Highlighting}
\end{Shaded}

\pandocbounded{\includegraphics[keepaspectratio]{analisis_ipc_files/analisis_ipc_files/figure-latex/unnamed-chunk-9-1.pdf}}

\begin{Shaded}
\begin{Highlighting}[]
\FunctionTok{ggsave}\NormalTok{(}\FunctionTok{here}\NormalTok{(}\StringTok{"salidas"}\NormalTok{, }\StringTok{"02\_divisiones\_julio.png"}\NormalTok{), grafico\_divisiones, }\AttributeTok{width =} \DecValTok{8}\NormalTok{, }\AttributeTok{height =} \DecValTok{5}\NormalTok{, }\AttributeTok{dpi =} \DecValTok{150}\NormalTok{)}
\end{Highlighting}
\end{Shaded}

\begin{Shaded}
\begin{Highlighting}[]
\CommentTok{\# Gráfico 3 (opcional): promedio de variación mensual por división, todo el periodo}
\NormalTok{promedio\_division }\OtherTok{\textless{}{-}}\NormalTok{ divisiones }\SpecialCharTok{\%\textgreater{}\%}
  \FunctionTok{group\_by}\NormalTok{(division) }\SpecialCharTok{\%\textgreater{}\%}
  \FunctionTok{summarise}\NormalTok{(}\AttributeTok{var\_mensual\_prom =} \FunctionTok{mean}\NormalTok{(var\_mensual, }\AttributeTok{na.rm =} \ConstantTok{TRUE}\NormalTok{)) }\SpecialCharTok{\%\textgreater{}\%}
  \FunctionTok{arrange}\NormalTok{(}\FunctionTok{desc}\NormalTok{(var\_mensual\_prom))}

\NormalTok{grafico\_promedio }\OtherTok{\textless{}{-}} \FunctionTok{ggplot}\NormalTok{(promedio\_division, }\FunctionTok{aes}\NormalTok{(}\AttributeTok{x =} \FunctionTok{reorder}\NormalTok{(division, var\_mensual\_prom), }\AttributeTok{y =}\NormalTok{ var\_mensual\_prom)) }\SpecialCharTok{+}
  \FunctionTok{geom\_col}\NormalTok{(}\AttributeTok{fill =} \StringTok{"\#E07A5F"}\NormalTok{) }\SpecialCharTok{+}
  \FunctionTok{coord\_flip}\NormalTok{() }\SpecialCharTok{+}
  \FunctionTok{labs}\NormalTok{(}
    \AttributeTok{title =} \StringTok{"Variación mensual promedio del IPC por división de gasto"}\NormalTok{,}
    \AttributeTok{subtitle =} \StringTok{"Promedio enero{-}julio 2026"}\NormalTok{,}
    \AttributeTok{x =} \ConstantTok{NULL}\NormalTok{, }\AttributeTok{y =} \StringTok{"Variación mensual promedio (\%)"}\NormalTok{,}
    \AttributeTok{caption =} \StringTok{"Fuente: DANE, boletines técnicos IPC ene{-}jul 2026 (anexos)"}
\NormalTok{  ) }\SpecialCharTok{+}
  \FunctionTok{theme\_minimal}\NormalTok{(}\AttributeTok{base\_size =} \DecValTok{12}\NormalTok{)}

\NormalTok{grafico\_promedio}
\end{Highlighting}
\end{Shaded}

\pandocbounded{\includegraphics[keepaspectratio]{analisis_ipc_files/analisis_ipc_files/figure-latex/unnamed-chunk-10-1.pdf}}

\begin{Shaded}
\begin{Highlighting}[]
\FunctionTok{ggsave}\NormalTok{(}\FunctionTok{here}\NormalTok{(}\StringTok{"salidas"}\NormalTok{, }\StringTok{"03\_promedio\_divisiones.png"}\NormalTok{), grafico\_promedio, }\AttributeTok{width =} \DecValTok{8}\NormalTok{, }\AttributeTok{height =} \DecValTok{5}\NormalTok{, }\AttributeTok{dpi =} \DecValTok{150}\NormalTok{)}
\end{Highlighting}
\end{Shaded}

```

\end{document}
